\chapter{Pulse Oximetry}

\section*{Overview}
Pulse oximetry is a quick, noninvasive method of estimating blood oxygen saturation and measuring pulse rate. It uses a small sensor, usually placed on a fingertip, toe, or earlobe.

\section*{Why It Is Used}
It is used to assess oxygenation in people with breathing or heart conditions, monitor patients during illness or procedures, guide oxygen therapy, and screen for low oxygen levels.

\section*{How It Is Performed}
The sensor shines light through the tissue and detects changes associated with oxygenated and deoxygenated hemoglobin. The device displays oxygen saturation (SpO2) and pulse rate within seconds.

\section*{Preparation and Limitations}
Remove dark nail polish when possible and keep the hand still and warm. Poor circulation, movement, skin pigmentation, nail products, cold extremities, and carbon monoxide exposure can affect accuracy. Readings should be interpreted with symptoms and the clinical situation.

\section*{Risks and Results}
Pulse oximetry is painless and generally has no significant risk. A persistently low reading may require clinical assessment, supplemental oxygen, or additional testing such as an arterial blood gas; the device alone does not diagnose the cause of low oxygen.

\section*{References}
\begin{enumerate}
  \item MedlinePlus. Pulse oximetry. U.S. National Library of Medicine; 2025.
  \item Current Medical Diagnosis and Treatment. McGraw-Hill; 2025.
\end{enumerate}
\clearpage
